% UC-4 Ordinamenti — Tabella delle Proprietà (versione ltablex corretta)
% Include con: % UC-4 Ordinamenti — Tabella delle Proprietà (versione ltablex corretta)
% Include con: % UC-4 Ordinamenti — Tabella delle Proprietà (versione ltablex corretta)
% Include con: % UC-4 Ordinamenti — Tabella delle Proprietà (versione ltablex corretta)
% Include con: \input{UC3 - Proprietà.txt}
% Preambolo richiesto:
% \usepackage{ltablex}
% \keepXColumns
\renewcommand{\arraystretch}{1.25}
\setlength{\tabcolsep}{6pt}

\begin{tabularx}{\textwidth}{|X|X|X|X|}
\caption{UC-3 Ordinamenti -- Proprietà e Regole di Validazione}
\label{tab:uc3-proprieta} \\ \hline
\textbf{Nome della Proprietà} & \textbf{Tipo} & \textbf{Regole di validazione} & \textbf{Note} \\ \hline
\endfirsthead

\hline
\textbf{Nome della Proprietà} & \textbf{Tipo} & \textbf{Regole di validazione} & \textbf{Note} \\ \hline
\endhead

\hline
\multicolumn{4}{r}{\small Continua nella pagina successiva} \\ \hline
\endfoot

\hline
\endlastfoot

\textbf{Colonne ordinabili} & String &
Obbligatorio, click sulla colonna per ordinarlo &
Definisce su quali campi è permesso applicare l'ordinamento (nome, \textit{stato}, \textit{priorità}, \textit{data creazione}) \\ \hline

\textbf{Direzione ordinamento} & Boolean (↑/↓) &
Invertibile dall'utente &
Indica se l'ordinamento è crescente o decrescente \\ \hline

\textbf{Indicatore visivo} & Icona UI &
Aggiornato automaticamente &
Mostra la direzione e il livello dell'ordinamento accanto all'intestazione \\ \hline

\textbf{Persistenza ordinamento} & URL param &
Facoltativa &
Se attiva, il sistema salva i criteri nella sessione o nell'URL \\ \hline

\textbf{Colonne non ordinabili} & String &
Click ignorato &
Colonna Tag \\ \hline

\end{tabularx}
% Preambolo richiesto:
% \usepackage{ltablex}
% \keepXColumns
\renewcommand{\arraystretch}{1.25}
\setlength{\tabcolsep}{6pt}

\begin{tabularx}{\textwidth}{|X|X|X|X|}
\caption{UC-3 Ordinamenti -- Proprietà e Regole di Validazione}
\label{tab:uc3-proprieta} \\ \hline
\textbf{Nome della Proprietà} & \textbf{Tipo} & \textbf{Regole di validazione} & \textbf{Note} \\ \hline
\endfirsthead

\hline
\textbf{Nome della Proprietà} & \textbf{Tipo} & \textbf{Regole di validazione} & \textbf{Note} \\ \hline
\endhead

\hline
\multicolumn{4}{r}{\small Continua nella pagina successiva} \\ \hline
\endfoot

\hline
\endlastfoot

\textbf{Colonne ordinabili} & String &
Obbligatorio, click sulla colonna per ordinarlo &
Definisce su quali campi è permesso applicare l'ordinamento (nome, \textit{stato}, \textit{priorità}, \textit{data creazione}) \\ \hline

\textbf{Direzione ordinamento} & Boolean (↑/↓) &
Invertibile dall'utente &
Indica se l'ordinamento è crescente o decrescente \\ \hline

\textbf{Indicatore visivo} & Icona UI &
Aggiornato automaticamente &
Mostra la direzione e il livello dell'ordinamento accanto all'intestazione \\ \hline

\textbf{Persistenza ordinamento} & URL param &
Facoltativa &
Se attiva, il sistema salva i criteri nella sessione o nell'URL \\ \hline

\textbf{Colonne non ordinabili} & String &
Click ignorato &
Colonna Tag \\ \hline

\end{tabularx}
% Preambolo richiesto:
% \usepackage{ltablex}
% \keepXColumns
\renewcommand{\arraystretch}{1.25}
\setlength{\tabcolsep}{6pt}

\begin{tabularx}{\textwidth}{|X|X|X|X|}
\caption{UC-3 Ordinamenti -- Proprietà e Regole di Validazione}
\label{tab:uc3-proprieta} \\ \hline
\textbf{Nome della Proprietà} & \textbf{Tipo} & \textbf{Regole di validazione} & \textbf{Note} \\ \hline
\endfirsthead

\hline
\textbf{Nome della Proprietà} & \textbf{Tipo} & \textbf{Regole di validazione} & \textbf{Note} \\ \hline
\endhead

\hline
\multicolumn{4}{r}{\small Continua nella pagina successiva} \\ \hline
\endfoot

\hline
\endlastfoot

\textbf{Colonne ordinabili} & String &
Obbligatorio, click sulla colonna per ordinarlo &
Definisce su quali campi è permesso applicare l'ordinamento (nome, \textit{stato}, \textit{priorità}, \textit{data creazione}) \\ \hline

\textbf{Direzione ordinamento} & Boolean (↑/↓) &
Invertibile dall'utente &
Indica se l'ordinamento è crescente o decrescente \\ \hline

\textbf{Indicatore visivo} & Icona UI &
Aggiornato automaticamente &
Mostra la direzione e il livello dell'ordinamento accanto all'intestazione \\ \hline

\textbf{Persistenza ordinamento} & URL param &
Facoltativa &
Se attiva, il sistema salva i criteri nella sessione o nell'URL \\ \hline

\textbf{Colonne non ordinabili} & String &
Click ignorato &
Colonna Tag \\ \hline

\end{tabularx}
% Preambolo richiesto:
% \usepackage{ltablex}
% \keepXColumns
\renewcommand{\arraystretch}{1.25}
\setlength{\tabcolsep}{6pt}

\begin{tabularx}{\textwidth}{|X|X|X|X|}
\caption{UC-2 Associzazione Etichette Personalizzate -- Proprietà e Regole di Validazione}
\label{tab:uc2-proprieta} \\ \hline
\textbf{Nome della Proprietà} & \textbf{Tipo} & \textbf{Regole di validazione} & \textbf{Note} \\ \hline
\endfirsthead

\hline
\textbf{Nome della Proprietà} & \textbf{Tipo} & \textbf{Regole di validazione} & \textbf{Note} \\ \hline
\endhead

\hline
\multicolumn{4}{r}{\small Continua nella pagina successiva} \\ \hline
\endfoot

\hline
\endlastfoot

\textbf{Nome} & String &
Obbligatorio, univoco, lunghezza da 1 a 30 caratteri. & L'aggiunta o la modifica aggiornano la descrizione della issue. \\ \hline

\textbf{Colore} & String (HEX) & Facoltativo, formato HEX (#RRGGBB), default:#808080. & Permette di personalizzare il colore della nuova etichetta, scegliendo tra un set di colori disponibili \hline

\textbf{Creatore} & String &
Generata automaticamente &
identificativo dell'Utente autenticato che la crea. \\ \hline

\end{tabularx}